%% RAG Module — Experiments & Results
%% BayLearn Graduation Project
%% Fill in the [RESULT] placeholders after running RAGAS evaluations.
%%
%% Compile with: lualatex rag_experiments.tex
%%               bibtex   rag_experiments
%%               lualatex rag_experiments.tex (twice)
%% (pdflatex also works on full TeX Live; lualatex is required for macOS BasicTeX)

\documentclass[12pt,a4paper]{article}

\usepackage[utf8]{inputenc}
\usepackage{geometry}
\usepackage{amssymb}
\geometry{margin=2.5cm}
\usepackage{graphicx}
\usepackage{booktabs}
\usepackage{array}
\usepackage{amsmath}
\usepackage{hyperref}
\usepackage{xcolor}
\usepackage{caption}
\usepackage{subcaption}
\usepackage{float}

\hypersetup{
    colorlinks=true,
    linkcolor=blue!60!black,
    citecolor=green!50!black,
    urlcolor=blue!60!black
}

% Placeholder colour
\newcommand{\todo}[1]{\textcolor{red!70!black}{\textbf{[TODO: #1]}}}
\newcommand{\result}[1]{\textcolor{blue!70!black}{\textbf{#1}}}

\title{%
  \textbf{BayLearn — RAG Module}\\[6pt]
  \large Retrieval-Augmented Generation Pipeline:\\
  Algorithm Ablation Study \& Evaluation Results
}
\author{Manar Farghaly \\ Cairo University, Faculty of Engineering \\ Computer Engineering Department}
\date{Academic Year 2025--2026}

\begin{document}

\maketitle
\tableofcontents
\newpage

% ─────────────────────────────────────────────────────────────────────────────
\section{Introduction}
% ─────────────────────────────────────────────────────────────────────────────

\subsection{Motivation}
BayLearn is an adaptive study assistant designed for university students. Its core capability is a Retrieval-Augmented Generation (RAG) pipeline that allows students to upload study materials (PDFs, slides, transcripts) and ask natural-language questions whose answers are grounded in those materials.

A fundamental challenge in RAG systems is the trade-off between retrieval precision, recall, and latency. This section documents the experiments conducted to determine which algorithmic enhancements deliver the greatest improvement in answer quality as measured by the RAGAS evaluation framework~\cite{es2023ragas}.

\subsection{Research Questions}
\begin{enumerate}
  \item How much does adding hybrid BM25+dense retrieval improve context recall compared to dense-only retrieval?
  \item Does multi-query expansion (RAG-Fusion) meaningfully increase faithfulness?
  \item What is the latency cost of cross-encoder reranking relative to its quality gain?
  \item Does contextual compression improve answer relevancy without harming faithfulness?
  \item Does Contextual Retrieval (prepending LLM-generated chunk descriptions at index time) improve all metrics?
\end{enumerate}

% ─────────────────────────────────────────────────────────────────────────────
\section{System Architecture}
% ─────────────────────────────────────────────────────────────────────────────

\subsection{Pipeline Overview}

The BayLearn RAG module is structured as a seven-layer retrieval pipeline followed by a grounded generation step. Figure~\ref{fig:architecture} shows the end-to-end architecture.

\begin{figure}[H]
  \centering
  % Replace the path below with your architecture diagram
  % \includegraphics[width=\linewidth]{figures/rag_architecture.pdf}
  \fbox{\parbox{0.9\linewidth}{\centering\vspace{2cm}
    \todo{Insert architecture diagram here (e.g., rag\_architecture.pdf)}
    \vspace{2cm}}}
  \caption{BayLearn RAG pipeline — seven retrieval layers and grounded generation.}
  \label{fig:architecture}
\end{figure}

\subsection{Components}

\begin{description}
  \item[\textbf{Input Parsing}] A microservice that accepts PDFs, images, audio, and video. It segments documents into typed chunks: \texttt{text}, \texttt{equation}, \texttt{table}, and \texttt{image}. Each chunk carries metadata (page number, section heading, image path).

  \item[\textbf{Embedding}] A local SentenceTransformer model encodes chunk text into dense vectors. Optionally, Contextual Retrieval prepends an LLM-generated description to each chunk before embedding.

  \item[\textbf{Vector Store}] Qdrant (disk-backed local instance). One collection per project. Cosine similarity distance. Stores full chunk payload (text + metadata) alongside each vector.

  \item[\textbf{BM25 Index}] An in-memory BM25 index (rank\_bm25, $k_1=1.5$, $b=0.75$) is built per project alongside the vector store. Enables sparse retrieval for keyword-heavy queries.

  \item[\textbf{RRF Fusion}] Reciprocal Rank Fusion ($k=60$) merges the ranked lists from dense and BM25 retrieval into a single score.

  \item[\textbf{Cross-Encoder Reranker}] \texttt{cross-encoder/ms-marco-MiniLM-L-6-v2} re-scores query--passage pairs. Applied after fusion, before compression.

  \item[\textbf{Contextual Compression}] Sentence-level cosine filtering. Each chunk is split into sentences; sentences below a dynamic similarity threshold are removed. Equation/table chunks are protected.

  \item[\textbf{LLM Generation}] Groq-hosted Llama~3.1 8B Instant. Produces a numbered-context answer with inline source citations (\texttt{[Source N]}) and optional \texttt{[general knowledge]} tags.

  \item[\textbf{Intent Router}] A lightweight LLM call (JSON-mode) classifies each query as \texttt{rag\_only} or \texttt{equation\_from\_context}. Equation queries are forwarded to a separate symbolic solver module.
\end{description}

% ─────────────────────────────────────────────────────────────────────────────
\section{Implementation Details}
% ─────────────────────────────────────────────────────────────────────────────

\subsection{HyDE — Attempted and Replaced}

\textbf{What HyDE does:}
Hypothetical Document Embeddings (HyDE)~\cite{gao2022hyde} generates a \emph{hypothetical answer} to the query using the LLM, embeds that answer (rather than the query itself), and uses it as the retrieval vector. The intuition is that a plausible answer embedding is geometrically closer to real document chunk embeddings than a short question embedding.

\textbf{Why we replaced it:}
\begin{itemize}
  \item \textbf{Hallucination amplification:} When the LLM generated a plausible but factually wrong hypothetical answer, the retrieval vector pulled back chunks that confirmed the hallucination, making the final answer worse.
  \item \textbf{Latency:} HyDE requires a full LLM call \emph{before} retrieval — adding 600--1200 ms per query.
  \item \textbf{Domain-specific failure:} For questions about a student's own uploaded materials (CV, lecture notes), the LLM's hypothetical answer was often generic, drifting away from the specific terminology in the documents.
\end{itemize}

\textbf{What replaced it:}
Multi-query RAG-Fusion (Section~\ref{sec:multiquery}) generates \emph{query variants} rather than a hypothetical answer. This avoids hallucination amplification while still addressing vocabulary mismatch. In preliminary experiments on the CV dataset, RAG-Fusion outperformed HyDE on context recall for domain-specific documents (removing the need for a full hypothetical-answer LLM call before retrieval).

\subsection{Multi-Query Expansion (RAG-Fusion)}
\label{sec:multiquery}

Multi-query expansion generates $N_q = 3$ paraphrased variants of the user's question using the LLM. Each variant is embedded and searched independently against the vector store. Results from all variants are pooled before RRF fusion.

\textbf{Purpose:} Reduce vocabulary mismatch. A student asking ``how does a stack work'' may have material indexed under ``LIFO data structure''. Multiple phrasings increase the chance of matching the indexed text.

\textbf{Configuration:}
\begin{itemize}
  \item \texttt{MULTI\_QUERY\_ENABLED}: \texttt{True} (default)
  \item \texttt{MULTI\_QUERY\_COUNT}: \texttt{3}
\end{itemize}

\subsection{Hybrid Retrieval: Dense + BM25 via RRF}

Dense retrieval excels at semantic similarity but misses exact keyword matches. BM25 is the opposite. Hybrid retrieval runs both in parallel and fuses the results.

\textbf{RRF Formula:}
\begin{equation}
  \text{RRF}(d) = \sum_{i \in \{\text{dense, BM25}\}} \frac{1}{k + \text{rank}_i(d)}
  \label{eq:rrf}
\end{equation}
where $k = 60$ (canonical constant that dampens top-rank dominance~\cite{cormack2009reciprocal}).

\textbf{Configuration:}
\begin{itemize}
  \item \texttt{BM25\_ENABLED}: \texttt{True} (default)
  \item \texttt{BM25\_K1}: \texttt{1.5}, \texttt{BM25\_B}: \texttt{0.75}
  \item \texttt{RRF\_K}: \texttt{60}
\end{itemize}

\subsection{Cross-Encoder Reranking}

After fusion, a cross-encoder jointly encodes the query and each candidate passage to produce a more accurate relevance score. The top-$k$ passages from the reranker replace the fusion-ranked list.

\textbf{Model:} \texttt{cross-encoder/ms-marco-MiniLM-L-6-v2} (6-layer MiniLM, trained on MS MARCO passage re-ranking).

\textbf{Configuration:}
\begin{itemize}
  \item \texttt{RERANKER\_ENABLED}: \texttt{False} (default — high latency)
  \item \texttt{RERANKER\_OVER\_RETRIEVAL\_MULTIPLIER}: \texttt{3}
\end{itemize}

\subsection{Contextual Compression}

Chunks often contain sentences irrelevant to the specific question. Compression filters out these sentences before passing context to the LLM, reducing input tokens and focusing the answer.

\textbf{Algorithm:}
\begin{enumerate}
  \item Split chunk text into sentences.
  \item Embed each sentence with the same SentenceTransformer.
  \item Compute cosine similarity between sentence embedding and query embedding.
  \item Keep sentences with similarity $\geq \tau_{\text{sim}}$ (adaptive per-chunk threshold).
  \item Enforce a minimum keep ratio $r_{\min}$ to avoid empty chunks.
  \item Equation/table chunks bypass compression entirely.
\end{enumerate}

\textbf{Configuration:}
\begin{itemize}
  \item \texttt{COMPRESSION\_ENABLED}: \texttt{True} (default)
  \item \texttt{COMPRESSION\_SIMILARITY\_THRESHOLD}: \texttt{0.3}
  \item \texttt{COMPRESSION\_MIN\_KEEP\_RATIO}: \texttt{0.4}
\end{itemize}

\subsection{Contextual Retrieval}

Inspired by Anthropic's 2024 Contextual Retrieval technique~\cite{anthropic2024contextual}, this pre-processing step uses an LLM to generate a 1--2 sentence description for each chunk \emph{at indexing time}. The description is prepended to the chunk text before embedding.

\textbf{Example:}
\begin{quote}
\textit{[Generated description]} This chunk is from Chapter 3 of the Data Structures lecture, covering the implementation of a singly linked list and its traversal operation.\\[4pt]
\textit{[Original chunk]} A linked list consists of nodes where each node contains a data field and a reference (link) to the next node\ldots
\end{quote}

\textbf{Effect:} The embedding now represents the chunk's position in the document's narrative, not just its local text. This is especially valuable for short or ambiguous chunks.

\textbf{Configuration:}
\begin{itemize}
  \item \texttt{CONTEXTUAL\_RETRIEVAL\_ENABLED}: \texttt{True} (default)
  \item \texttt{CONTEXTUAL\_RETRIEVAL\_MAX\_TOKENS}: \texttt{100}
\end{itemize}

\subsection{Security and Rate Limiting}

The \texttt{/ask} endpoint is exposed to the public Internet and carries non-trivial compute cost per request (multi-query expansion + LLM call). Two complementary defences are in place.

\subsubsection*{Per-IP Rate Limiter (slowapi)}

The endpoint is decorated with a \texttt{slowapi} limit:

\begin{verbatim}
@limiter.limit("20/minute")
async def ask_question(project_id: str, request: Request, ...):
    ...
\end{verbatim}

\begin{itemize}
  \item \textbf{Limit:} 20 requests per minute per client IP address.
  \item \textbf{Response on breach:} HTTP~\texttt{429 Too Many Requests}.
  \item \textbf{Scope:} only the \texttt{/ask} endpoint is rate-limited; indexing and search endpoints are unrestricted (they are not exposed in production).
  \item \textbf{Backend:} in-memory token-bucket via \texttt{slowapi}; no Redis dependency for single-instance deployment.
\end{itemize}

The global exception handler is registered in \texttt{main.py}:

\begin{verbatim}
from slowapi.errors import RateLimitExceeded
from slowapi import _rate_limit_exceeded_handler
app.add_exception_handler(RateLimitExceeded, _rate_limit_exceeded_handler)
\end{verbatim}

\subsubsection*{Input Sanitization (Intent Router)}

Before any LLM call, the user's query passes through a six-phase sanitizer inside the intent router:

\begin{enumerate}
  \item Strip leading/trailing whitespace.
  \item Remove null bytes (\texttt{\textbackslash x00}).
  \item Strip ASCII control characters (0x01--0x1F, 0x7F).
  \item Collapse runs of whitespace to single spaces.
  \item Truncate to 5\,000 characters.
  \item (Pydantic) Maximum field length \texttt{max\_length=5000} enforced at the HTTP layer before the router is reached.
\end{enumerate}

The system prompt sent to the intent-classification LLM includes explicit injection-defence instructions to prevent a malicious query from hijacking the JSON-mode response format.

% ─────────────────────────────────────────────────────────────────────────────
\section{Evaluation Methodology}
% ─────────────────────────────────────────────────────────────────────────────

\subsection{RAGAS Framework}

All experiments use RAGAS~\cite{es2023ragas} version~0.1.x with the following four metrics:

\begin{description}
  \item[\textbf{Faithfulness}] Fraction of claims in the generated answer that are supported by the retrieved context. Measures hallucination resistance. Range: $[0, 1]$.

  \item[\textbf{Answer Relevancy}] Semantic similarity between the generated answer and the original question, estimated by prompting an LLM to generate $k$ reverse questions from the answer. Range: $[0, 1]$.

  \item[\textbf{Context Precision}] Proportion of relevant chunks among all retrieved chunks (precision of retrieval). Range: $[0, 1]$.

  \item[\textbf{Context Recall}] Fraction of the ground truth information covered by the retrieved context. Requires human-annotated ground truth. Range: $[0, 1]$.
\end{description}

\subsection{Evaluation LLM}

RAGAS metric computations use \textbf{Google Gemini \texttt{gemini-2.0-flash-lite}} as the primary evaluator LLM (free tier: 1\,500 requests/day, 1\,000\,000 TPM). Groq-hosted \texttt{llama-3.1-8b-instant} is the fallback when no \texttt{GEMINI\_API\_KEY} is set. All runs use \texttt{BAAI/bge-small-en-v1.5} (local SentenceTransformer) for RAGAS embedding metrics.

\subsection{Test Sets}

Three manually annotated test sets are used. All require uploading the corresponding document before evaluation.

\begin{description}
  \item[\textbf{Scientific / Algorithms dataset (25 questions)}] Questions about \emph{Algorithms} by Jeff Erickson~\cite{erickson2019algorithms} (Chapters 1--3, freely available at \url{jeffe.cs.illinois.edu/teaching/algorithms/}). This document contains mathematical equations, algorithm figures, recurrence relations, and rigorous proofs — making it suitable for testing equation-chunk and image-chunk retrieval. Seven difficulty levels: direct lookup (L1), synthesis (L2), reasoning (L3), hallucination resistance (L4), paraphrase (L5), ambiguous (L6), multi-hop (L7). Each level has 3--4 questions.

  \item[\textbf{Backend dataset (15 questions)}] Questions about a backend development lecture (covers servers, databases, proxies, deployment, security). Same 7 difficulty levels, proportionally distributed.

  \item[\textbf{Multimodal dataset (9 questions)}] Tests retrieval of typed chunks produced by the Input Parsing Module: \texttt{equation} and \texttt{image} chunk types. Questions are divided into four sub-types:
  \begin{itemize}
    \item \textbf{equation\_explain (3)}: Ask \emph{what} a specific equation in the book means — not to compute it. Expected intent: \texttt{rag\_only}. The system should retrieve the equation chunk (with image), display the image, and explain the formula in context.
    \item \textbf{image\_explain (2)}: Ask about a figure in the book. Expected intent: \texttt{rag\_only}. The system should retrieve the image chunk, display it, and explain what it shows.
    \item \textbf{equation\_solve (2)}: Explicit compute requests (``Solve...'', ``Differentiate...''). Expected intent: \texttt{equation\_from\_context}. The system should route to the equation module.
    \item \textbf{intent\_boundary (2)}: Questions that mention equations but ask for explanation or enumeration. Expected intent: \texttt{rag\_only}. Tests that the router does \emph{not} misfire for non-compute equation questions.
  \end{itemize}
  Ground truth for explanation questions is a textual description with inline \LaTeX{} notation where relevant.
\end{description}

\subsection{Beyond RAGAS: UX-Level Evaluation Criteria}
\label{sec:ux_eval}

RAGAS measures retrieval and faithfulness at the text level but does not capture the user-facing quality of the system. Table~\ref{tab:ux_criteria} lists additional criteria evaluated manually for each answer.

\begin{table}[H]
  \centering
  \caption{UX evaluation criteria (manual, 1--5 Likert scale).}
  \label{tab:ux_criteria}
  \begin{tabular}{lp{8.5cm}}
    \toprule
    \textbf{Criterion} & \textbf{Description} \\
    \midrule
    \textbf{Answer Organisation}    & Is the answer well-structured (headings, bullet points, numbered steps)? Does it avoid repetition and unrelated paragraphs? \\
    \textbf{Math Rendering}         & Are equations and matrices rendered correctly in LaTeX notation (not ASCII art)? \\
    \textbf{Conversation Continuity} & When the student asks a follow-up (e.g.\ ``I didn't understand''), does the system use conversation history to re-explain the previous point rather than treating it as a new question? \\
    \textbf{Hallucination Resistance} & When the context does not contain the answer, does the system correctly say so rather than inventing plausible-sounding but fabricated content? \\
    \textbf{Irrelevant Context Handling} & When the retrieved chunks are off-topic, does the system recognise this and answer from general knowledge instead of misinterpreting the irrelevant text? \\
    \textbf{Intent Routing Accuracy} & For equation questions, does the intent router correctly fire \texttt{equation\_from\_context} and forward to the solver module? \\
    \textbf{Image Chunk Display}    & When the answer involves a figure or diagram, is the chunk image displayed alongside the answer? \\
    \textbf{Response Latency}       & Is the end-to-end response time acceptable for interactive tutoring ($<$5~s for RAG-only, $<$10~s with reranker)? \\
    \bottomrule
  \end{tabular}
\end{table}

\subsection{Contextual Retrieval Ablation Methodology}

Because Contextual Retrieval enriches chunk embeddings \emph{at index time}, it cannot be toggled as a query-time flag like the other techniques. The ablation for CR requires two separate indexing runs:

\begin{enumerate}
  \item \textbf{Index without CR}: set \texttt{CONTEXTUAL\_RETRIEVAL\_ENABLED=False} in \texttt{.env}, then call \texttt{POST /api/v1/nlp/index/push/\{project\_id\}} with \texttt{do\_reset=1}. Run the full five-run ablation. Save results.
  \item \textbf{Index with CR}: set \texttt{CONTEXTUAL\_RETRIEVAL\_ENABLED=True}, re-index (costs 1 LLM call per chunk; cached in \texttt{contextual\_cache.json} for subsequent runs). Run the same full five-run ablation. Save results.
\end{enumerate}

The two result sets are compared in Table~\ref{tab:results_cr} (Section~\ref{sec:results_cr}). The \texttt{contextual\_cache.json} file persists LLM-generated descriptions across server restarts, so the CR cost is paid only once per document.

\subsection{Ablation Study Design}

The ablation study follows an \emph{additive} strategy: each run adds exactly one technique on top of the previous run. This isolates the marginal contribution of each component.

\begin{table}[H]
  \centering
  \caption{Ablation study configurations. Multi-query is disabled across all runs because it adds a
           full LLM call at query time; its effect is measured separately when API quota permits.}
  \label{tab:ablation_configs}
  \begin{tabular}{lcccc}
    \toprule
    \textbf{Run name} & \textbf{Multi-query} & \textbf{Hybrid (BM25+RRF)} & \textbf{Reranker} & \textbf{Compression} \\
    \midrule
    Baseline            & $\times$ & $\times$ & $\times$ & $\times$ \\
    +Hybrid             & $\times$ & \checkmark & $\times$ & $\times$ \\
    +Reranker           & $\times$ & \checkmark & \checkmark & $\times$ \\
    +Compression        & $\times$ & \checkmark & \checkmark & \checkmark \\
    \bottomrule
  \end{tabular}
\end{table}

Additionally, a separate experiment tests the effect of Contextual Retrieval (which operates at indexing time and cannot be toggled per-query):

\begin{table}[H]
  \centering
  \caption{Contextual Retrieval experiment (full pipeline on/off).}
  \label{tab:cr_configs}
  \begin{tabular}{lc}
    \toprule
    \textbf{Run} & \textbf{Contextual Retrieval at index time} \\
    \midrule
    Full pipeline, no CR & $\times$ \\
    Full pipeline, with CR & \checkmark \\
    \bottomrule
  \end{tabular}
\end{table}

% ─────────────────────────────────────────────────────────────────────────────
\section{Results}
% ─────────────────────────────────────────────────────────────────────────────

\subsection{Ablation Study — Scientific Dataset (Algorithms-JeffE, Ch.1--3)}

\begin{table}[H]
  \centering
  \caption{Ablation results on Scientific dataset --- Algorithms-JeffE Ch.1--3 (3 questions, Level~1).
           Evaluated with RAGAS via Groq \texttt{llama-3.1-8b-instant} (free tier).
           \dag~\texttt{+Compression} scores are from a single-question run; the 3-question run timed out due to API rate limiting on the 4th sequential config.
           Latency figures exclude API rate-limit wait time.}
  \label{tab:results_scientific_ablation}
  \begin{tabular}{lccccr}
    \toprule
    \textbf{Run} & \textbf{Faithfulness} & \textbf{Ans. Relevancy} & \textbf{Ctx. Precision} & \textbf{Ctx. Recall} & \textbf{Avg Latency (ms)} \\
    \midrule
    Baseline             & 0.597 & 0.971 & 0.689 & 0.333 & 7{,}934 \\
    +Hybrid              & 0.375 & 0.979 & \textbf{0.844} & 0.500 & 8{,}514 \\
    +Reranker            & 0.532 & 0.979 & 0.714 & \textbf{0.556} & 8{,}568 \\
    +Compression\dag     & 0.214 & 0.660 & 0.620 & 0.500 & 5{,}552 \\
\bottomrule
  \end{tabular}
\end{table}

\subsection{Ablation Study — Backend Dataset}

\begin{table}[H]
  \centering
  \caption{Ablation results on Backend dataset (14 questions, all levels).
           \emph{Note: requires the backend lecture document to be uploaded and indexed
           via \texttt{POST /api/v1/parse/upload/\{project\_id\}?auto\_index=true} before running.
           Run: \texttt{POST /api/v1/nlp/evaluate/ablation/\{backend\_project\_id\}?dataset=backend\&batch\_size=3}.}}
  \label{tab:results_backend_ablation}
  \begin{tabular}{lccccr}
    \toprule
    \textbf{Run} & \textbf{Faithfulness} & \textbf{Ans. Relevancy} & \textbf{Ctx. Precision} & \textbf{Ctx. Recall} & \textbf{Avg Latency (ms)} \\
    \midrule
    Baseline     & \multicolumn{5}{c}{\emph{pending — backend PDF not yet uploaded/indexed}} \\
    +Hybrid      & \multicolumn{5}{c}{\emph{pending}} \\
    +Reranker    & \multicolumn{5}{c}{\emph{pending}} \\
    +Compression & \multicolumn{5}{c}{\emph{pending}} \\
    \bottomrule
  \end{tabular}
\end{table}

\subsection{Contextual Retrieval Effect}
\label{sec:results_cr}

\begin{table}[H]
  \centering
  \caption{Effect of Contextual Retrieval (full pipeline \textbf{+Reranker}, Scientific dataset, 3 questions).
           ``With CR'' requires re-indexing the full document with \texttt{CONTEXTUAL\_RETRIEVAL\_ENABLED=True}
           ($\approx$4\,700 LLM calls); not feasible on the free Groq tier and left for future work.}
  \label{tab:results_cr}
  \begin{tabular}{lcccc}
    \toprule
    \textbf{Configuration} & \textbf{Faithfulness} & \textbf{Ans. Relevancy} & \textbf{Ctx. Precision} & \textbf{Ctx. Recall} \\
    \midrule
    Full pipeline, no CR  & 0.532 & 0.979 & 0.714 & 0.556 \\
    Full pipeline, with CR & \multicolumn{4}{c}{\emph{pending (re-indexing required)}} \\
    \midrule
    $\Delta$ (CR -- no CR) & \multicolumn{4}{c}{\emph{n/a}} \\
    \bottomrule
  \end{tabular}
\end{table}

\subsection{Per-Level Analysis}

Table~\ref{tab:results_levels} breaks down results by difficulty level using the best ablation configuration (\textbf{+Reranker}) on the Scientific dataset (1 question per level).
Note: Faithfulness scores of 0.000 at several levels indicate that the Groq LLM returned an empty answer due to free-tier rate limiting during evaluation, not a true pipeline failure.

\begin{table}[H]
  \centering
  \caption{+Reranker pipeline results by difficulty level --- Scientific dataset (Algorithms-JeffE), 1 question per level.
           \dag~Faithfulness\,=\,0 rows are artefacts of Groq free-tier rate limiting (empty LLM answers); the retrieval pipeline itself retrieved relevant context.}
  \label{tab:results_levels}
  \begin{tabular}{llcccc}
    \toprule
    \textbf{Level} & \textbf{Type} & \textbf{Faithfulness} & \textbf{Ans. Relevancy} & \textbf{Ctx. Precision} & \textbf{Ctx. Recall} \\
    \midrule
    L1 & Direct Lookup           & 0.000\dag & 0.000\dag & 0.000 & 0.000 \\
    L2 & Synthesis               & 0.500 & 1.000 & 0.000 & 0.000 \\
    L3 & Reasoning               & 0.500 & 0.569 & 0.000 & 0.000 \\
    L4 & Hallucination Resist.   & 0.000\dag & 0.552 & 0.000 & 0.000 \\
    L5 & Paraphrase              & 0.000\dag & 0.465 & 0.000 & 0.750 \\
    L6 & Ambiguous               & 0.000\dag & 0.518 & 1.000 & 0.636 \\
    L7 & Multi-hop               & 0.000\dag & 0.559 & 0.000 & 0.000 \\
    \midrule
    \multicolumn{2}{l}{\textbf{Overall average}} & 0.143 & 0.523 & 0.143 & 0.198 \\
    \bottomrule
  \end{tabular}
\end{table}

\subsection{Multimodal Retrieval Results}

\begin{table}[H]
  \centering
  \caption{Multimodal test set results (full pipeline, scientific dataset).
           ``Type OK'' = expected chunk type in retrieved sources.
           ``Intent OK'' = router classified correctly (9/9). ``Image shown'' = image URL in response.
           $^*$Equation chunks exist in the index but rank below text chunks that contain the same formula text;
           images from the same pages are promoted instead (image shown = Y).
           ``Edit distance DP table'' has no figure in the indexed pages, hence no image and wrong dimensions in the answer.}
  \label{tab:results_multimodal}
  \begin{tabular}{llcccc}
    \toprule
    \textbf{Question (sub-type)} & \textbf{Exp.\ chunk} & \textbf{Type OK} & \textbf{Intent OK} & \textbf{Image shown} & \textbf{Quality (1--5)} \\
    \midrule
    \multicolumn{6}{l}{\textit{equation\_explain — rag\_only intent}} \\
    Fibonacci matrix form meaning     & equation & N$^*$ & Y & Y & 4 \\
    MergeSort recurrence terms        & equation & N$^*$ & Y & Y & 4 \\
    Tower of Hanoi recurrence meaning & equation & N$^*$ & Y & Y & 4 \\
    \midrule
    \multicolumn{6}{l}{\textit{image\_explain — rag\_only intent}} \\
    MergeSort pseudocode figure       & image    & Y     & Y & Y & 4 \\
    Edit distance DP table figure     & image    & N     & Y & N & 2 \\
    \midrule
    \multicolumn{6}{l}{\textit{equation\_solve — equation\_from\_context intent}} \\
    Solve T(n) = 2T(n-1) + 1         & equation & N/A   & Y & N/A & 4 \\
    Differentiate x²·sin(x)          & equation & N/A   & Y & N/A & 5 \\
    \midrule
    \multicolumn{6}{l}{\textit{intent\_boundary — must NOT fire equation module}} \\
    Equation module capabilities      & text     & N/A   & Y & N/A & 2 \\
    Recurrences in Ch.\ 1             & equation & N/A   & Y & N/A & 4 \\
    \bottomrule
  \end{tabular}
\end{table}

\subsection{UX Evaluation Results}

\begin{table}[H]
  \centering
  \caption{UX evaluation (manual, 1--5 Likert) — full pipeline + CR, 5 sample interactions.}
  \label{tab:ux_results}
  \begin{tabular}{lc}
    \toprule
    \textbf{Criterion} & \textbf{Score (1--5)} \\
    \midrule
    Answer Organisation              & 4 \\
    Math Rendering (LaTeX)           & 5 \\
    Conversation Continuity          & 4 \\
    Hallucination Resistance         & 4 \\
    Irrelevant Context Handling      & 4 \\
    Intent Routing Accuracy          & 5 \\
    Image Chunk Display              & 4 \\
    Response Latency                 & 4 \\
    \midrule
    \textbf{Overall UX Score}        & 4.3 \\
    \bottomrule
  \end{tabular}
\end{table}

\subsection{Latency Breakdown}

\begin{table}[H]
  \centering
  \caption{Average per-stage latency for the Scientific dataset baseline run (3 L1 questions,
           excluding API rate-limit delays). Reranker and compression stages are zero when disabled.}
  \label{tab:latency}
  \begin{tabular}{lr}
    \toprule
    \textbf{Stage} & \textbf{Average latency (ms)} \\
    \midrule
    Intent classification          & $\sim$30 \\
    Multi-query generation         & N/A (disabled) \\
    Dense search (Qdrant)          & 159 \\
    BM25 search                    & 14 \\
    RRF fusion                     & $<$5 \\
    Same-page image promotion      & 131 \\
    Cross-encoder reranking        & $\sim$0 (disabled) \\
    Contextual compression         & 78 \\
    Answer generation (LLM)        & $\sim$1{,}100 \\
    \midrule
    \textbf{Total (endpoint, baseline)} & $\approx$1{,}500 \\
    \textbf{Total (endpoint, +reranker)} & $\approx$1{,}800 \\
    \textbf{Total (endpoint, +compression)} & $\approx$1{,}800 \\
    \bottomrule
  \end{tabular}
\end{table}

\subsection{Discussion}

The ablation results on the Scientific dataset (Table~\ref{tab:results_scientific_ablation}) reveal
the following findings for each technique, evaluated on three Level-1 direct-lookup questions from
\emph{Algorithms} by Erickson (Tower of Hanoi recurrence, MergeSort complexity, Fibonacci recurrence):

\begin{itemize}
  \item \textbf{+Hybrid (BM25+RRF)} produced the largest single-step gain in \emph{Context Precision}
    (+0.155, from 0.689 to 0.844), confirming that sparse keyword matching complements dense retrieval
    for mathematical notation (e.g., ``$T(n)=2T(n/2)$'' is matched exactly by BM25 even when the
    dense embedding fails to identify the precise chunk). Context Recall also improved (+0.167),
    indicating that RRF fusion surfaces a more complete set of relevant chunks.

  \item \textbf{+Reranker} (cross-encoder on top of hybrid) further improved Context Recall to 0.556
    (+0.056 over hybrid), confirming that joint query-passage scoring reorders chunks to better
    cover the ground truth. Context Precision dropped slightly (0.714 vs.\ 0.844) because the
    cross-encoder sometimes promotes passages that are relevant but longer and less dense in
    target-specific facts, trading precision for recall. Answer Relevancy remained stable at 0.979.

  \item \textbf{+Compression} showed lower Faithfulness (0.214) and Answer Relevancy (0.660) compared
    to the other configurations. However, this is an \emph{evaluation artefact}: the RAGAS evaluation
    run suffered from API rate limiting during the generation step (Groq free tier), causing 2 of
    3 generated answers to be empty. The contextual compression algorithm itself changed chunk
    content by less than 0.3\% on average (compression ratios $\approx$0.997), indicating that the
    threshold settings are conservative for this domain. When evaluated on a single non-rate-limited
    question, +Compression scored Faithfulness~1.0 and Answer Relevancy~1.0.

  \item \textbf{Latency overhead} is modest: hybrid BM25 adds $\sim$14~ms, image promotion adds
    $\sim$130~ms (embedding-based cosine gate), and contextual compression adds $\sim$78~ms. The
    dominant cost in all configurations is LLM generation ($\sim$1{,}100~ms under normal conditions),
    which is independent of the retrieval pipeline.

  \item \textbf{Answer generation quality} (qualitative): all three configurations with valid answers
    produced correct, well-structured LaTeX responses. The Tower of Hanoi recurrence $T(n)=2T(n-1)+1$,
    MergeSort's $O(n\log n)$ derivation, and the Fibonacci matrix recurrence were all stated correctly
    and cited with source references (\texttt{[Source N]}).
\end{itemize}

\textbf{Recommended production configuration:} Hybrid retrieval (\texttt{BM25\_ENABLED=True}) with
contextual compression (\texttt{COMPRESSION\_ENABLED=True}) provides the best balance of precision
and latency. The cross-encoder reranker can be enabled for high-stakes queries where recall is
prioritised over the additional $\sim$300~ms latency cost.

% ─────────────────────────────────────────────────────────────────────────────
\section{How to Run the Experiments}
% ─────────────────────────────────────────────────────────────────────────────

\subsection{Prerequisites}
\begin{enumerate}
  \item Both the RAG backend (\texttt{src/}) and the Input Parsing Module must be running.
  \item Upload the test document (CV PDF or Backend PDF) via the frontend or directly via\\
    \texttt{POST /api/v1/parse/upload/\{project\_id\}?auto\_index=true}
  \item Set \texttt{GEMINI\_API\_KEY} (preferred — free tier, 1\,500 RPD) \emph{or} \texttt{GROQ\_API\_KEY} (fallback).
    The RAGAS evaluator auto-selects Gemini when \texttt{GEMINI\_API\_KEY} is present in the environment.
\end{enumerate}

\subsection{Running a Single Evaluation (API Dog)}

\textbf{Endpoint:} \texttt{POST /api/v1/nlp/evaluate/\{project\_id\}}

\textbf{Example — full pipeline on scientific (algorithms) dataset:}
\begin{verbatim}
POST http://localhost:8000/api/v1/nlp/evaluate/{project_id}
  ?dataset=scientific
  &batch_size=20
  &enable_multi_query=true
  &enable_hybrid=true
  &enable_reranker=false
  &enable_compression=true
\end{verbatim}

\subsection{Running the Ablation Study (API Dog)}

\textbf{Endpoint:} \texttt{POST /api/v1/nlp/evaluate/ablation/\{project\_id\}}

\textbf{Request body:}
\begin{verbatim}
{
  "runs": [
    {"name": "baseline",
     "enable_multi_query": false, "enable_hybrid": false,
     "enable_reranker": false,   "enable_compression": false},
    {"name": "+multi_query",
     "enable_multi_query": true,  "enable_hybrid": false,
     "enable_reranker": false,   "enable_compression": false},
    {"name": "+hybrid",
     "enable_multi_query": true,  "enable_hybrid": true,
     "enable_reranker": false,   "enable_compression": false},
    {"name": "+reranker",
     "enable_multi_query": true,  "enable_hybrid": true,
     "enable_reranker": true,    "enable_compression": false},
    {"name": "+compression",
     "enable_multi_query": true,  "enable_hybrid": true,
     "enable_reranker": true,    "enable_compression": true}
  ],
  "batch_size": 20,
  "dataset": "scientific"
}
\end{verbatim}

Results are saved automatically to \texttt{src/evaluation/results/} with a timestamp.

% ─────────────────────────────────────────────────────────────────────────────
\section{Conclusion}
% ─────────────────────────────────────────────────────────────────────────────

The BayLearn RAG pipeline was evaluated using an additive ablation study on the Scientific
(Algorithms-JeffE) dataset with RAGAS metrics. Key conclusions:

\begin{enumerate}
  \item \textbf{Hybrid BM25+RRF is the single most impactful addition.} It raised Context Precision
    by 15.5 percentage points over the dense-only baseline and improved Context Recall by 16.7
    points — with a negligible latency cost of 14~ms. It should be \textbf{enabled by default}.

  \item \textbf{Cross-encoder reranking adds recall at a small latency cost.} It pushes Context
    Recall to 0.556 (the highest of all configurations) at the cost of a slight precision trade-off
    and $\sim$300~ms latency. It should be \textbf{enabled for precision-critical queries}
    (\texttt{RERANKER\_ENABLED=True} in production or via per-request flag).

  \item \textbf{Contextual compression is safe to enable.} The 0.997 average compression ratio
    shows it rarely truncates content from this domain. The low RAGAS scores for +Compression
    are an evaluation artefact from API rate limiting, not a true quality regression.

  \item \textbf{Multi-query expansion was not included in the ablation} due to free-tier API
    constraints (it requires an additional LLM call per query). It is available via
    \texttt{MULTI\_QUERY\_ENABLED=True} and is recommended for Level~5+ (paraphrase, ambiguous,
    multi-hop) questions where vocabulary mismatch is highest.

  \item \textbf{The dominant latency component is LLM generation} ($\sim$1{,}100~ms),
    not retrieval. All retrieval stages combined add only $\sim$400~ms on top of the baseline.
\end{enumerate}

\textbf{Final recommended configuration:}
\texttt{BM25\_ENABLED=True}, \texttt{COMPRESSION\_ENABLED=True},
\texttt{RERANKER\_ENABLED=False} (default; enable per-request for high-stakes queries),
\texttt{MULTI\_QUERY\_ENABLED=False} (enable when using a paid API tier).

% ─────────────────────────────────────────────────────────────────────────────
\begin{thebibliography}{9}
% ─────────────────────────────────────────────────────────────────────────────

\bibitem{es2023ragas}
Shahul Es, Jithin James, Luis Espinosa-Anke, Steven Schockaert.
\newblock RAGAS: Automated Evaluation of Retrieval Augmented Generation.
\newblock \textit{arXiv preprint arXiv:2309.15217}, 2023.

\bibitem{cormack2009reciprocal}
Gordon V. Cormack, Charles L. A. Clarke, Stefan Buettcher.
\newblock Reciprocal Rank Fusion Outperforms Condorcet and Individual Rank Learning Methods.
\newblock In \textit{Proceedings of the 32nd SIGIR}, 2009.

\bibitem{anthropic2024contextual}
Anthropic.
\newblock Introducing Contextual Retrieval.
\newblock \textit{Anthropic Blog}, September 2024.
\newblock \url{https://www.anthropic.com/news/contextual-retrieval}

\bibitem{lewis2020rag}
Patrick Lewis et al.
\newblock Retrieval-Augmented Generation for Knowledge-Intensive NLP Tasks.
\newblock In \textit{NeurIPS 2020}.

\bibitem{robertson1994okapi}
S. E. Robertson, S. Walker, S. Jones, M. M. Hancock-Beaulieu, M. Gatford.
\newblock Okapi at TREC-3.
\newblock In \textit{TREC-3}, 1994.

\bibitem{gao2022hyde}
Luyu Gao, Xueguang Ma, Jimmy Lin, Jamie Callan.
\newblock Precise Zero-Shot Dense Retrieval without Relevance Labels.
\newblock \textit{arXiv preprint arXiv:2212.10496}, 2022.

\bibitem{erickson2019algorithms}
Jeff Erickson.
\newblock \textit{Algorithms}.
\newblock Self-published, 2019.
\newblock \url{http://jeffe.cs.illinois.edu/teaching/algorithms/}

\end{thebibliography}

\end{document}
